% Source: https://artofproblemsolving.com/community/c3h2651057_geo_question_i_had
% Stable identity: aops-topic-2651057-q01
The slopes are
\[m_{PQ}=\dfrac{9-6}{-7+5}=-\dfrac32,\qquad m_{QR}=\dfrac{10-9}{1+7}=\dfrac18,\qquad m_{PR}=\dfrac{10-6}{1+5}=\dfrac23.\]
Since $m_{PQ}\cdot m_{PR}=-\dfrac32\cdot\dfrac23=-1$, sides $PQ$ and $PR$ are perpendicular: the triangle has a right angle at $P$. Its legs are $PQ=\sqrt{2^2+3^2}=\sqrt{13}$ and $PR=\sqrt{6^2+4^2}=2\sqrt{13}$, so the area is $\dfrac12\cdot\sqrt{13}\cdot2\sqrt{13}=\boxed{13}$.

So the slopes are $\boxed{-\dfrac32,\ \dfrac18,\ \dfrac23}$, and the triangle is right-angled at $P$.

\begin{asy}
usepackage("amsmath");
size(5cm);
defaultpen(fontsize(10pt));
import graph;
pair P=(-5,6),Q=(-7,9),R=(1,10);
draw((-8.5,0)--(2.5,0),gray,Arrow(4));draw((0,-0.5)--(0,11.5),gray,Arrow(4));
for(int i=-8;i<=2;++i) draw((i,-0.12)--(i,0.12),gray);
fill(P--Q--R--cycle,lightgray);draw(P--Q--R--cycle,linewidth(0.8));
draw(P+0.5*unit(Q-P)--P+0.5*unit(Q-P)+0.5*unit(R-P)--P+0.5*unit(R-P));
dot(P);dot(Q);dot(R);
label("$P(-5,6)$",P,S);label("$Q(-7,9)$",Q,W);label("$R(1,10)$",R,NE);
label("$-\dfrac32$",(P+Q)/2+(-0.3,0),W,fontsize(8pt));label("$\dfrac18$",(Q+R)/2,N,fontsize(8pt));label("$\dfrac23$",(P+R)/2,SE,fontsize(8pt));
\end{asy}
